\section{SF-E848W: a quartic complete intersection outside the earlier explicit Artin tests}\label{SF-E848W:section}

Author: /root/sixfolds, GPT-6 Astra Pro; inherited effort not exposed.
Date: 2026-11-22. Author, root and independent whole-proof readings passed without theorem repair. The frozen mathematical original is retained.
This canonical conversion has a separate full-text check.
This is a comparison theorem for the auxiliary RM fourfold programme.
It is not an auxiliary-to-target Weil-eightfold transfer theorem.

The exact arithmetic supplement has passed, with its initial checker-sign
failure preserved. A full written proof follows. No actual solution of the
first-order companion equations is asserted. The scope of the comparison,
including what it does not establish about strictness, is stated in Section 7.

\subsection{1. The fixed algebra, bivectors, and assertions}\label{SF-E848W:s1}

Let \(R=\mathbf C[[x,y,z,w]]\), and put

\[
\begin{split}
 f_1&=x^4+xy^3,& f_2&=y^4-8x^3y,\\
 f_3&=z^4+zw^3,& f_4&=w^4-8z^3w,\\
 I&=(f_1,f_2,f_3,f_4),& A&=R/I,\\
 \pi_+&=\partial_x\wedge\partial_z,&
 \pi_-&=\partial_y\wedge\partial_w.
\end{split}
\tag{QW.1}
\]

The two bivectors have complementary two-dimensional supports. All formal
coordinate comparisons below transport both the ideal and these tensors;
they do not change the tensors to a different geometric pair.
Let \(P_+\) and \(P_-\) denote the six bracket entries \(\{f_i,f_j\}_+\) and \(\{f_i,f_j\}_-\)
in \(A\), with \(i<j\). For an alternating four-by-four matrix use the convention

\[
 \operatorname{Pf}(P)=P_{12}P_{34}-P_{13}P_{24}+P_{14}P_{23}.
\tag{QW.2}
\]

The following properties hold.

\begin{enumerate}
\def\labelenumi{\arabic{enumi}.}
\tightlist
\item
  \(I\) is a complete intersection of colength 256, with embedding dimension
  four, relation order four, and one-dimensional socle in degree twelve.
\item
  All scalar cross-minors \(P_{+,e}P_{-,f}-P_{-,e}P_{+,f}\) vanish. Both pure
  bracket maps have rank four, their combined image has dimension eight,
  and every mixed pencil has nonzero Pfaffian.
\item
  For every minimal formal generating set of \(I\) and every bracket pencil,
  no two elements in its entry ideal form a nonzero socle pair of E842:
  if their product is nonzero, their principal-ideal intersection has
  dimension at least five, rather than one.
\item
  \(I\) has no formal presentation by powers of a regular parameter system.
  No three independent minimal generators commute under either of the
  two specified bivectors, even after arbitrary formal generator changes.
\item
  Every common-CI matrix extension over this \(A\) is an actual local module
  covered by E848. It can be realized in a coherent balanced-character
  fourfold example using the already constructed balanced hull. These
  examples are excluded by E848; none is claimed to be semiregular.
\end{enumerate}

The universal statement in item 3 is the substantive comparison with the
socle-pair test. A single failed choice of entries would not suffice.

\subsection{2. A binary quartic algebra and its exact multiplication}\label{SF-E848W:s2}

Set

\[
 B=\mathbf C[x,y]/(x^4+xy^3,\ y^4-8x^3y),
 \qquad \tau_B=x^3y^3.
\tag{QW.3}
\]

The two homogeneous polynomials have no common projective zero. If one
coordinate vanishes, the other fourth power forces the other coordinate
to vanish. If both are nonzero, the equations give \(x^3=-y^3\) and \(y^3=8x^3\),
which would give \(9x^3=0\). Thus they are coprime. Unique factorization shows
that the second is a nonzerodivisor modulo the first: divisibility of its
product with a polynomial by the first implies divisibility of that
polynomial by the first. Hence the two-term Koszul resolution is exact,
and the Hilbert series of \(B\) is

\[
 \frac{(1-t^4)^2}{(1-t)^2}
       =(1+t+t^2+t^3)^2.
\tag{QW.4}
\]

The sixteen monomials \(x^iy^j\) with \(0\le i,j\le3\) span. To see this directly,
use \(x^4=-xy^3\) and \(y^4=8x^3y\). A monomial with \(i\ge4\) and \(j\ge1\) returns to
itself with coefficient -8 after the two reductions, so it is zero.
The same argument applies with \(i\ge1\) and \(j\ge4\). A pure power of degree
four, five, or six reduces into the sixteen displayed monomials; a pure
power of higher degree becomes one of the vanishing mixed monomials.
The Hilbert-series dimension is sixteen, so the spanning monomials are
a basis. In particular \(\tau_B\) is nonzero, and \(B\) has graded dimensions
\((1,2,3,4,3,2,1)\).

The characteristic-zero complete-intersection Jacobian-socle theorem,
in exactly the form retained in E842.4, identifies the socle. Here its
normalization is also explicit:

\[
 \det\begin{pmatrix}
 4x^3+y^3&3xy^2\\
 -24x^2y&4y^3-8x^3
 \end{pmatrix}=144\tau_B\quad\hbox{in }B.
\tag{QW.5}
\]

For example \(x^6=-\tau_B\) and \(y^6=8\tau_B\), and substitution into the
determinant gives \((80+32+32)\tau_B\). Thus the nonzero top class is the
socle generator. The invoked source is Eisenbud--Ulrich, J. reine
angew. Math. 756 (2019), Theorem 7.1 p.~210, with Appendix A.1/A.3/A.5,
pp.~222--225.\cite{SFEisenbudUlrich2019} This is the already completely checked source input of
E836--E844, not a newly claimed whole-source reading.

Write the two derivative columns as

\[
 \begin{split}
 u_1&=4x^3+y^3,&u_2&=-24x^2y,\\
 v_1&=3xy^2,&v_2&=-8x^3+4y^3,\\
 U&=\langle u_1,u_2\rangle,&V&=\langle v_1,v_2\rangle.
 \end{split}
\tag{QW.6}
\]

In the ordered basis \((x^3,x^2y,xy^2,y^3)\) of \(B_3\), the four columns
\(u_1,u_2,v_1,v_2\) have determinant -1728. Consequently \(U\) and \(V\) are
complementary two-dimensional subspaces of \(B_3\). Direct reduction gives

\[
 (u_i v_j)_{i,j=1,2}
       =72\tau_B\begin{pmatrix}0&1\\-1&0\end{pmatrix}.
\tag{QW.7}
\]

For completeness, \(u_1v_1\) and \(u_2v_2\) are sums of the vanishing mixed
monomials above. The product \(u_1v_2\) is \(-32x^6+8x^3y^3+4y^6=72\tau_B\),
and \(u_2v_1=-72\tau_B\). These signs will be used in both cross-minor and
Pfaffian calculations.

\subsection{3. Every nonzero linear form is injective on both derivative planes}\label{SF-E848W:s3}

Let \(L=\alpha x+\beta y\) be a nonzero linear form. In the ordered basis
\((x^3y,x^2y^2,xy^3)\) of \(B_4\), multiplication by \(L\) on \(U\) and \(V\) has matrices

\[
 M_U(L)=\begin{pmatrix}
 12\beta&-24\alpha\\
 0&-24\beta\\
 -3\alpha&0
 \end{pmatrix},\qquad
 M_V(L)=\begin{pmatrix}
 0&24\beta\\
 3\alpha&0\\
 3\beta&12\alpha
 \end{pmatrix}.
\tag{QW.8}
\]

The first and second row minor of \(M_U\) is \(-288\beta^2\), and its first
and third row minor is \(-72\alpha^2\). The second and third row minor of
\(M_V\) is \(36\alpha^2\), and its first and third row minor is \(-72\beta^2\).
At least one listed minor for each matrix is nonzero whenever \(L\) is
nonzero. Both multiplication maps are therefore injective. This is
a statement for every complex linear form, not just coordinate forms
or a finite sampled set.

The completed algebra \(A\) in (QW.1) is the tensor product of two copies
of \(B\), with variables \((x,y)\) and \((z,w)\); completion does not change a
finite-dimensional quotient. The regular sequence follows either by
the tensor product of the two exact Koszul resolutions or by doing
the second coprime pair over the first free finite-dimensional factor.
Thus \(A\) has length 256 and Hilbert function

\[
 (\dim A_d)_{d=0}^{12}
   =(1,4,10,20,31,40,44,40,31,20,10,4,1).
\tag{QW.9}
\]

Its socle is generated by \(\tau=\tau_B\otimes\tau_B=x^3y^3z^3w^3\),
and the four-variable Jacobian is \(20736\tau\). Introduce the subspace

\[
 W=(U\otimes U)\oplus(V\otimes V)\subset A_6.
\tag{QW.10}
\]

It has dimension eight because \(U\) and \(V\) are complementary in each
binary factor.

We claim that multiplication by every nonzero \(L\) in \(A_1\) is injective
on \(W\). Write \(L=L_1+L_2\) according to the two binary factors. For
\(p=p_U+p_V\) in \(W\), the two terms \(L_1p\) and \(L_2p\) have distinct bidegrees
\((4,3)\) and \((3,4)\), so they cannot cancel. If \(L_1\) is nonzero, then
\(L_1p_U\) and \(L_1p_V\) lie respectively in \(B_4\otimes U\) and \(B_4\otimes V\).
These are disjoint subspaces. The injectivity just proved for each
binary derivative plane implies \(p_U=p_V=0\). If \(L_1=0\), use the same
argument in the second factor, where \(L_2\) is nonzero. This proves

\[
 0\ne p\in W\quad\Longrightarrow\quad
       (A_1\xrightarrow{\,p\,}A_7)\text{ is injective}.
\tag{QW.11}
\]

The reversal of the quantified variables in this displayed conclusion
is valid: a nonzero linear form killing a nonzero \(p\) would contradict
the preceding injectivity on \(W\).

\subsection{4. Bracket ranks, all cross-minors, and the mixed Pfaffian}\label{SF-E848W:s4}

Use \(u_i\) and \(v_i\) also for the identical expressions in the second
binary factor. The only nonzero bracket positions are

\[
 P_{+,i,2+j}=u_i\otimes u_j,\qquad
 P_{-,i,2+j}=v_i\otimes v_j\quad(1\le i,j\le2),
 \qquad P_{\pm,12}=P_{\pm,34}=0.
\tag{QW.12}
\]

The four entries of each pure tuple are independent; their spans are
\(U\otimes U\) and \(V\otimes V\). The two bracket maps from \(\bigwedge^2\mathbf C^4\)
therefore have rank four and kernel dimension two. Their combined
span is \(W\). Every nonzero pencil \(aP_++bP_-\) also has rank four: if
both coefficients are nonzero, project a linear relation onto the
two direct summands in (QW.10); if one coefficient is zero, use the
pure rank statement.

Let \(\epsilon_{12}=1\), \(\epsilon_{21}=-1\), and \(\epsilon_{11}=\epsilon_{22}=0\).
For \(e=(i,2+j)\) and \(f=(k,2+l)\), equation (QW.7) gives

\[
 P_{+,e}P_{-,f}=72^2\epsilon_{ik}\epsilon_{jl}\tau
            =72^2\epsilon_{ki}\epsilon_{lj}\tau
            =P_{-,e}P_{+,f}.
\tag{QW.13}
\]

If an index is 12 or 34 both sides are zero. Thus every scalar
cross-minor vanishes, including all 36 ordered choices of \(e,f\).
The pure Pfaffians vanish because each original bivector is
decomposable, or by direct multiplication in (QW.12). The same
calculation with (QW.7) gives the precise mixed value

\[
 \operatorname{Pf}(aP_++bP_-)=-20736ab\,\tau.
\tag{QW.14}
\]

It is nonzero when \(ab\) is nonzero. Its sign agrees with the chosen
coordinate order and (QW.2): \(\pi_+\wedge\pi_-\) reverses \(y\) and \(z\).
Consequently the vanished scalar cross-minors do not imply a zero
Jacobian socle for this actual complete intersection.

\subsection{5. No socle pair in any presentation or any pencil entry ideal}\label{SF-E848W:s5}

We first record the elementary graded duality implication needed
here. In a finite graded algebra with one-dimensional socle in degree
\(s\), multiplication \(A_j\otimes A_{s-j}\) to \(A_s\) is a perfect pairing.
Indeed the principal ideal of any nonzero homogeneous element meets
the socle: choose a nonzero product of it with maximal possible
further degree. Every element of the maximal ideal kills
that maximal product, so it lies in the socle. Its homogeneous
degree is \(s\). This proves nondegeneracy in both variables and hence
perfectness. The argument applies to our socle identified above.

Apply this pairing with \(s=12\) to (QW.11). The transpose of
multiplication \(p:A_1\) to \(A_7\) is multiplication \(p:A_5\) to \(A_{11}\).
Therefore

\[
            pA_5=A_{11}\qquad(0\ne p\in W).
\tag{QW.15}
\]

Let \(J\) be the ideal generated by all entries of \(P_+\) and \(P_-\). Every
element of \(J\) has order at least six. Its degree-six initial form,
if nonzero, belongs to \(W\). Let \(a\) in \(J\) have nonzero initial degree-six
form \(a_6\). From (QW.15), the initial forms of \(aA_5\) span \(A_{11}\). The
degree-twelve piece \(A_{12}\) is contained in \(aA\) as well. To see this,
first choose a degree-eleven initial value \(z_{11}\) of a product \(ar_5\)
with a linear form \(L\) for which \(Lz_{11}\) is nonzero. Such a choice
exists because \(A_1A_{11}=A_{12}\) is nonzero; for example a top monomial
is a variable times a degree-eleven monomial. Multiplication by \(L\)
kills the degree-twelve error in \(ar_5\), since \(A_{13}=0\), and gives
a nonzero element of \(A_{12}\) in \(aA\). Subtracting suitable multiples of
this top class from the lifts of each element of \(A_{11}\) now proves

\[
                  \mathfrak m_A^{11}=A_{11}\oplus A_{12}
                               \subseteq aA.
\tag{QW.16}
\]

This subspace has dimension five. If \(a,b\) in \(J\) have nonzero product,
then both have order six: each order is at least six, and a larger
order for either would force the product into \(A_{\ge13}=0\). Thus
(QW.16) applies to both, giving

\[
 0\ne ab\quad\Longrightarrow\quad
 \dim_{\mathbf C}(aA\cap bA)\ge5,
 \qquad aA\cap bA\ne\mathbf C\,ab.
\tag{QW.17}
\]

In particular such a pair cannot satisfy the hypothesis of E842a.
Its nonzero product is automatically in the degree-twelve socle,
so the failed hypothesis is exactly the one-dimensional intersection.

Now let \(g_1,\ldots,g_4\) be any minimal formal generating set of \(I\).
There is a matrix \(T\) in \(\operatorname{GL}_4(R)\) such that \(g=Tf\). Modulo \(I\), all terms
involving derivatives of \(T\) have a factor \(f_i\) and vanish. Hence

\[
       P_\pm(g)=\overline T\,P_\pm(f)\,\overline T^{\mathsf t}.
\tag{QW.18}
\]

The joint entry ideal \(J\) is unchanged because \(T\) is invertible, and
the entry ideal of any pencil is contained in \(J\). Equations
(QW.16)--(QW.17) therefore exclude socle pairs in every one of these
ideals, for every minimal presentation and every pencil. Formal
coordinate isomorphisms preserve principal-ideal intersections and
the socle. This establishes the claimed coordinate-independent
comparison for the fixed geometric pair of tensors.

\subsection{6. No power presentation and no commuting triple}\label{SF-E848W:s6}

Suppose \(I\) were generated by powers \(u_i^{a_i}\), where the \(u_i\) are a
formal regular parameter system and the \(a_i\) are positive integers.
Since every nonzero element of \(I\) has order at least four, every
\(a_i\) is at least four. The degree-four relation subspace has dimension
four. Terms with \(a_i>4\) cannot contribute to that subspace, while the
fourth powers of the independent linear initial parameters are
linearly independent. Thus all four exponents would have to equal
four, and the degree-four relation space would contain a nonzero
fourth power of a linear form.

It contains no such power. All coefficients of \(x_i^2x_j^2\), for
distinct variables, vanish on the span of \(f_1,\ldots,f_4\). For a fourth
power \(L^4\), these coefficients are \(6c_i^2c_j^2\), where the \(c_i\) are the
coefficients of \(L\). Over \(\mathbf C\) they can all vanish only if at most one
\(c_i\) is nonzero. Thus \(L^4\) would have to be a multiple of one
coordinate fourth power. None lies in the relation space: eliminating
the mixed monomials \(xy^3,x^3y,zw^3,z^3w\) forces all four relation
coefficients to be zero. This contradiction rules out every exponent
choice, permutation, and nonlinear formal parameter change.

For a commuting triple, let three independent members of a minimal
generating system have zero pairwise brackets under \(\pi_+\). Their
degree-four initial relation classes span a three-dimensional
subspace \(E\) of \(\mathbf C^4\). Taking degree-six initial brackets says that
\(\bigwedge^2E\), of dimension three, is contained in the kernel of the
pure bracket map. That kernel has dimension two by Section 4,
which is impossible. The same proof works for \(\pi_-\) and for every
formal minimal presentation, using (QW.18) in initial degree.
There are no linear relations to which the linear-generator part
of E836 could apply.

The relation order is four. Thus the square-algebra and regular
quadratic-initial-form hypotheses of E837--E840 are absent. The
homogeneous quadratic local-module hypothesis of E846 is absent
as well. These assertions concern the hypotheses of the named
results, not all conceivable consequences of the underlying
companion equations.

\subsection{7. Actual coherent realizations and the precise comparison scope}\label{SF-E848W:s7}

This is an actual algebra at a smooth point of a complex abelian
fourfold. Choose formal logarithm coordinates compatible with the
RM tangent splitting, so that \(\pi_+\) and \(\pi_-\) have the form (QW.1).
The ideal contains the thirteenth power of the maximal ideal.
The finite algebra therefore comes from the finite jet of the
algebraic local ring, so it defines an actual finite subscheme \(Z\).
No convergent or global coordinate system is being asserted.

For every \(q\ge1\) and arbitrary six matrices \(b\) in \(\operatorname{Mat}_q(A)\), the
Koszul group \(\operatorname{Ext}^1_R(I^q,A^q)=\operatorname{Mat}_q(A)^6\) gives an actual extension.
Its global realization is exactly E848 Section 9: take \(q\) copies
of the actual rank-32 balanced hull \(H_0\) from E815/E820, choose a
surjection \(H_0^{\oplus q}\) to \(\mathcal O_Z^q\), and extend its kernel by \(\mathcal O_Z^q\)
using the desired matrix class. The Ext spectral sequence has
finite support and realizes every local class. The completed
module is \(M_b\) direct-sum \(R^{31q}\). Its two finite defects have the
same character and cancel, so its entire character is that of
the balanced hull. This coherent sheaf has no negative self-Ext.
E848 excludes weak semiregularity for every specified effective
coherent translation/tensor action on it; the identity is one
actual action. No nontrivial action is asserted to exist.

The comparison now has an exact and limited content. For this
common CI, none of the reduced, curvilinear, linear-generator,
commuting-triple, parameter-power, regular-quadratic-jet, or
socle-pair hypotheses just listed applies; all scalar cross-minors
also vanish. E848 nevertheless applies uniformly to arbitrary
matrix extension tuples, including noncommuting tuples. The scalar
nonzero case was already excluded by the immediately preceding
E847; this example is not a new exclusion beyond that scalar
theorem. The result proves a strict separation from the named
earlier Artin tests, not existence of matrices \(B_e,C_+,C_-\) solving
all actual first-order equations. In particular it does not claim
that every tuple \(b\) survives a direct tuple-specific first-order
calculation. It does not claim a semiregular seed, nor any transfer
from the auxiliary fourfold to a target Weil eightfold.

\subsection{8. Exact supplement and preserved failed attempts}\label{SF-E848W:s8}

The separate script strict\_ci\_search/quartic\_comparison\_certificate.py
checks the displayed defining relations, all 4096 binary basis
associativity identities, the full Hilbert function and Jacobian,
all 36 cross-minors, both pure and the mixed Pfaffian, ranks 4/4/8,
the complementary derivative determinant, the universal two-by-two
minors of (QW.8), the four binary pairings, and the elementary
pure-power coefficient tests. It uses exact arithmetic. The
universal no-annihilator and all-presentation arguments above are
proofs, not extrapolations from these finite checks.

The successful run began at 03:20:16.929796 UTC and completed at
03:20:17.024641 UTC on 2026-11-22; internal elapsed time was
0.09485129799577408 seconds. It recorded 66 named checks. An initial
checker expected \(-12\alpha\) instead of \(+12\alpha\) in \(M_V\) and the
corresponding wrong sign of its alpha-square minor. Direct reduction
raised an assertion before success. The first source and the exact
failure/repair record are preserved alongside the final script.
The algebra and the injectivity argument did not change.

Earlier cubic searches are retained as failed comparisons: the
generic cubic candidate had nonzero scalar cross-minors, and four
corrected structured cubic candidates had injective cross-minor
maps. They were not strict witnesses. No independent reviewer or
new numerical execution is claimed by this paragraph.
